\section{KDD Governance Model}
\label{sec:kdd-governance}

The \kdd~governance model defines the rules by which agents and human contributors
create, modify, and consume artefacts across the repository ecosystem.

\subsection{Core Principles}

\begin{enumerate}
  \item \textbf{Knowledge first.} Every implementation decision must be grounded
    in a documented knowledge artefact: an \adr, a metric definition, a workflow
    specification, or a dataset card.
  \item \textbf{Traceability.} Every output artefact must carry metadata linking
    it to the decision, dataset, and workflow that produced it.
  \item \textbf{Reproducibility by default.} No result may be included in the
    paper without a corresponding generation script and dataset.
  \item \textbf{Separation of concerns.} Each repository has a single declared
    responsibility; cross-cutting concerns are mediated by interfaces.
\end{enumerate}

\subsection{Structured Document Definitions}

The \sdd~mechanism enforces document structure at generation time. When an agent
produces a new artefact, the \kdd~governance layer validates the artefact against
the registered schema for its type before committing it.

\subsection{Governance Metrics}

We define four primary governance metrics:

\begin{itemize}
  \item \metric{Traceability Score}: fraction of artefacts with complete metadata.
  \item \metric{SDD Compliance Rate}: fraction of structured documents passing
    schema validation.
  \item \metric{Audit Coverage}: fraction of tool calls recorded in the audit log.
  \item \metric{Evidence Density}: mean number of evidence citations per
    recommendation.
\end{itemize}
